\documentclass{article}
\usepackage{amsmath}
\usepackage[hidelinks]{hyperref}
\title{Common receive windows: where static rate ceilings suffice}
\author{}
\date{}
\begin{document}
\maketitle
\paragraph{Audit.}
The original live-writer result holds at its fixed radio settings,
but a stronger static comparator wins in its quiet transfer
window. Narrowing \emph{both} matching writers admits a first
copy and mechanically refuses a second before it can spend
receive budget. This relies on an A-before-B service envelope
and the availability of the narrower cap to both senders. It
does not distinguish identities and is not robust to all lawful
orderings at the original receiver budget. The decisive
comparison is not exact checks versus no checks. It is semantically
blind resource admission versus authorized elimination before the
resource actually at risk.

\paragraph{Sources and contract.}
The experiment rebuilds the exact pinned native SubT dispatch and
visibility model as in the parent; the selected cave projection
has the same two physical terminal samples and the same final
breadcrumb positions. Delivery of $P$'s literal current nine-byte
packet depends on the established first-representative RF result.
Independent fallback $Q$ has a declared live/drop RF seam. Native
dispatch charges bytes using the sending node's limit for its
own TX and for the terminal's current RX window. Traffic $X$ at
the terminal is actually emitted and can create background RX
debit. The separate native side flow continues throughout.
The two owed keys $A$ (nine bytes) and $B$ (ten bytes) share the
native receiving resource. Order and available copy opportunities
are part of the declared source envelope, not secretly selectable
by one policy only. Copies are submitted within a single native
window, with no refunds. The holdout cave/mission has not been
used. The test does not assert a stock interface for changing a
live radio limit or for pre-enqueue authorization.

\paragraph{Quiet-window comparator.}
For both orders with an A opportunity before B, both $A$ sender
caps at 17 and $R$ at 19 meet the two duties whenever at least
one A path is viable. The first reachable nine fits; another
would raise RX to 18, which the second writer's 17 refuses. B
fits at 19. TX can still be spent by the rejected duplicate,
including a physically failed attempt. Fresh exact-key lookup
and shared control of every old writer are not necessary in
this restricted envelope.

More generally suppose initial live RX is in $[L,H]$ and the
same sender may start with at most $U$ live TX bytes. If that
sender can be the sole first copy at both maxima, and can
genuinely reach the second RX test at the minimum background,
its admissible static cap satisfies
\[
                 \max(H+a,U+a)\le q<L+2a.
\]
For B after exactly one A, its RX threshold must be at least
$H+a+b$ (as well as any independent TX requirement). The bound
assumes the credits remain live at the tested cuts, compatible
first/second opportunity support, and an $a$-byte copy. Do not
apply it to a different background/timing law by carrying the
aggregate interval alone. The native sweep includes background
8 and 9. For example, if backgrounds from 0 through 9 are
possible, the bound that passes a first copy after 9 conflicts
with the bound that rejects a duplicate after 0.

\paragraph{Two fixed-budget discriminators.}
Hold the B cap at 19. Add lawful schedules where B precedes the
only viable A, with both choices of the surviving A path. Any
passing fixed pair must then allow 10+9 on each sender. But at
those caps both nine-byte copies can arrive before B in the
dual-path source history, leaving 18 and blocking B. All 1,089
configured cap pairs in 0--32 fail that mixed support. The
scoped exact guard at 19/19/19 passes, since it only suppresses
a later matching copy after an actual receipt; an unavailable
first transmitter never supplies that proof. Raising B's cap
to tolerate extra bytes is a changed resource alternative, not
ruled out by this claim.

For a separate discriminator let a distinct nine-byte $C$ share
$P$'s radio and be owed after B. If a different history needs Q
as the first viable A, its cap must admit nine. Serving C after
A and B requires $P$'s cap at least 28. With these fixed
thresholds the native Q-before-P dual path admits both copies
of A, exhausting B's 19 even if C itself eventually fits.
The current keyed guard suppresses $P$'s redundant A but lets
its distinct C through with exactly the same caps and actual
nominal P path. The tested dispatch substitute therefore lies
at a different effective resource boundary, not at an
application deduplication point downstream of an already incurred
credit charge. Established
\href{https://www.rfc-editor.org/rfc/rfc8938.html}{DetNet packet
elimination and forwarding-layer protection} remain valid
comparators if placed before the affected resource or with
genuinely separate reservations. The local exact-key guard itself
is not an advancement over equally provisioned elimination with
matching key/generation semantics and actual placement.

\paragraph{Verification.}
The independent byte/TX reference agrees with all 5078 native
rows, including narrow and wide resource variants, physical
first-route outcomes, distinct C and an unchanged separate
work stream. Lean checks the bound and distinct/mixed-order
arithmetic. The selected mode of the guard in this fixture has
serial visibility of actual current receipt and authorized
control of all pending matching writers. No claim is made about
multi-process atomicity, source support outside the explicit
orders, or a stronger site mission result.
\end{document}
